\documentclass[10pt,a4paper]{amsart}
\usepackage{amsmath,amssymb,amsthm,graphicx}
\usepackage{verbatim}
\newenvironment{reference}{\comment}{\endcomment}
%\newenvironment{reference}{}{}
\newenvironment{slogan}{\comment}{\endcomment}
\newenvironment{history}{\comment}{\endcomment}
% For commutative diagrams we use Xy-pic
\usepackage[all]{xy}

% We use 2cell for 2-commutative diagrams.
\xyoption{2cell}
% Theorem environments.
%
\theoremstyle{plain}
\newtheorem{theorem}[subsection]{Theorem}
\newtheorem{proposition}[subsection]{Proposition}
\newtheorem{lemma}[subsection]{Lemma}

\theoremstyle{definition}
\newtheorem{definition}[subsection]{Definition}
\newtheorem{example}[subsection]{Example}
\newtheorem{exercise}[subsection]{Exercise}
\newtheorem{situation}[subsection]{Situation}

\theoremstyle{remark}
\newtheorem{remark}[subsection]{Remark}
\newtheorem{remarks}[subsection]{Remarks}

\numberwithin{equation}{subsection}

% Macros
%
\def\lim{\mathop{\mathrm{lim}}\nolimits}
\def\colim{\mathop{\mathrm{colim}}\nolimits}
\def\Spec{\mathop{\mathrm{Spec}}}
\def\Hom{\mathop{\mathrm{Hom}}\nolimits}
\def\Ext{\mathop{\mathrm{Ext}}\nolimits}
\def\SheafHom{\mathop{\mathcal{H}\!\mathit{om}}\nolimits}
\def\SheafExt{\mathop{\mathcal{E}\!\mathit{xt}}\nolimits}
\def\Sch{\mathit{Sch}}
\def\Mor{\mathop{\mathrm{Mor}}\nolimits}
\def\Ob{\mathop{\mathrm{Ob}}\nolimits}
\def\Sh{\mathop{\mathit{Sh}}\nolimits}
\def\NL{\mathop{N\!L}\nolimits}
\def\CH{\mathop{\mathrm{CH}}\nolimits}
\def\proetale{{pro\text{-}\acute{e}tale}}
\def\etale{{\acute{e}tale}}
\def\QCoh{\mathit{QCoh}}
\def\Ker{\mathop{\mathrm{Ker}}}
\def\Im{\mathop{\mathrm{Im}}}
\def\Coker{\mathop{\mathrm{Coker}}}
\def\Coim{\mathop{\mathrm{Coim}}}

% Boxtimes
%
\DeclareMathSymbol{\boxtimes}{\mathbin}{AMSa}{"02}

%
% Macros for moduli stacks/spaces
%
\def\QCohstack{\mathcal{QC}\!\mathit{oh}}
\def\Cohstack{\mathcal{C}\!\mathit{oh}}
\def\Spacesstack{\mathcal{S}\!\mathit{paces}}
\def\Quotfunctor{\mathrm{Quot}}
\def\Hilbfunctor{\mathrm{Hilb}}
\def\Curvesstack{\mathcal{C}\!\mathit{urves}}
\def\Polarizedstack{\mathcal{P}\!\mathit{olarized}}
\def\Complexesstack{\mathcal{C}\!\mathit{omplexes}}
% \Pic is the operator that assigns to X its picard group, usage \Pic(X)
% \Picardstack_{X/B} denotes the Picard stack of X over B
% \Picardfunctor_{X/B} denotes the Picard functor of X over B
\def\Pic{\mathop{\mathrm{Pic}}\nolimits}
\def\Picardstack{\mathcal{P}\!\mathit{ic}}
\def\Picardfunctor{\mathrm{Pic}}
\def\Deformationcategory{\mathcal{D}\!\mathit{ef}}
\usepackage{cleveref}
\crefname{originalreference}{original reference}{original references}
\Crefname{originalreference}{Original reference}{Original references}
\makeatletter
\newlabel{item-intersection-quasi-compact-open}{{item-intersection-quasi-compact-open}{source}{}{sourcepointer-1}{}}
\newlabel{item-intersection-quasi-compact-open@cref}{{[originalreference][0][]item-intersection-quasi-compact-open}{[1][0][]source}{}{}{}}
\makeatother
\crefname{originallemma}{original lemma}{original lemmas}
\Crefname{originallemma}{Original lemma}{Original lemmas}
\makeatletter
\newlabel{lemma-directed-inverse-limit-spectral-spaces}{{lemma-directed-inverse-limit-spectral-spaces}{source}{}{sourcepointer-2}{}}
\newlabel{lemma-directed-inverse-limit-spectral-spaces@cref}{{[originallemma][0][]lemma-directed-inverse-limit-spectral-spaces}{[1][0][]source}{}{}{}}
\makeatother
\crefname{originallemma}{original lemma}{original lemmas}
\Crefname{originallemma}{Original lemma}{Original lemmas}
\makeatletter
\newlabel{lemma-make-spectral-space}{{lemma-make-spectral-space}{source}{}{sourcepointer-3}{}}
\newlabel{lemma-make-spectral-space@cref}{{[originallemma][0][]lemma-make-spectral-space}{[1][0][]source}{}{}{}}
\makeatother
\crefname{originallemma}{original lemma}{original lemmas}
\Crefname{originallemma}{Original lemma}{Original lemmas}
\makeatletter
\newlabel{lemma-spectral-sub}{{lemma-spectral-sub}{source}{}{sourcepointer-4}{}}
\newlabel{lemma-spectral-sub@cref}{{[originallemma][0][]lemma-spectral-sub}{[1][0][]source}{}{}{}}
\makeatother
\title{The points specializing to a constructible subset, minus the subset, form a spectral space}
\author{The Stacks Project Authors}
\date{stacks-project:topology.tex:local-snapshot-20261012; GNU FDL 1.2 or later}
\begin{document}
\maketitle
\noindent\textit{Editorial note.} All delivered mathematical text is the author's own, copied from the pinned Stacks Project chapter source sources/stacks-topology/pinned-original-source.tex; no statement or proof step was rewritten, reordered or removed.\par\medskip
\noindent\textit{Editorial note.} Preamble provenance: the Stacks Project keeps one single shared preamble for every chapter; this card takes it directly from the staged sources/stacks-topology/preamble.tex as three verbatim spans (lines 8-12, 14-18 and 154-217), so no derived file was written and nothing was added or redefined. The wrapper class and the base packages are supplied by the assembler.\par\medskip
\noindent\textit{Editorial note.} Bibliography provenance: neither the counted statement nor its proof contains a citation (this lane rescans the whole span before writing the card), so the card carries no bibliography and no entry was reconstructed.\par\medskip
\noindent\textit{Editorial note.} Selection provenance: this card was opened under the file-based control contract (CONTROL.md), with the mechanical pre-card gate run on its statement and proof interval before the card was written.\par\medskip
\noindent\textit{Allowed standard prerequisites (editorial).}\begin{itemize}
\item The description of a spectral subset as a limit of the quasi-compact opens containing it. Invoked at the start of the counted proof (source lines 5297-5298). Identifies the set of points specializing to the constructible subset as such a limit.
\item The statement that a subset closed in the constructible topology is spectral. Invoked inside the counted proof (source lines 5300-5302). Makes the differences inside the quasi-compact opens spectral.
\item The statement that an inverse limit of spectral spaces is spectral. Invoked at the end of the counted proof (source lines 5304-5305). Concludes spectrality of the limit.
\end{itemize}\medskip
\ifdefined\hypertarget\hypertarget{sourcepointer-1}{}\fi
\noindent\textit{Source reference.} Original reference item-intersection-quasi-compact-open: sources/stacks-topology/pinned-original-source.tex, source lines 5232--5232. Navigation pointer only: the counted text refers to the Stacks item item-intersection-quasi-compact-open, which is an existing labeled result of this same chapter and is not part of this card. The pointer carries the target's own label as its number and its definition line as its span; its content is not reproduced here.\par\medskip
\ifdefined\hypertarget\hypertarget{sourcepointer-2}{}\fi
\noindent\textit{Source reference.} Original lemma lemma-directed-inverse-limit-spectral-spaces: sources/stacks-topology/pinned-original-source.tex, source lines 5119--5119. Navigation pointer only: the counted text refers to the Stacks item lemma-directed-inverse-limit-spectral-spaces, which is an existing labeled result of this same chapter and is not part of this card. The pointer carries the target's own label as its number and its definition line as its span; its content is not reproduced here.\par\medskip
\ifdefined\hypertarget\hypertarget{sourcepointer-3}{}\fi
\noindent\textit{Source reference.} Original lemma lemma-make-spectral-space: sources/stacks-topology/pinned-original-source.tex, source lines 5222--5222. Navigation pointer only: the counted text refers to the Stacks item lemma-make-spectral-space, which is an existing labeled result of this same chapter and is not part of this card. The pointer carries the target's own label as its number and its definition line as its span; its content is not reproduced here.\par\medskip
\ifdefined\hypertarget\hypertarget{sourcepointer-4}{}\fi
\noindent\textit{Source reference.} Original lemma lemma-spectral-sub: sources/stacks-topology/pinned-original-source.tex, source lines 4645--4645. Navigation pointer only: the counted text refers to the Stacks item lemma-spectral-sub, which is an existing labeled result of this same chapter and is not part of this card. The pointer carries the target's own label as its number and its definition line as its span; its content is not reproduced here.\par\medskip
\begin{lemma}
\label{lemma-make-spectral-space-minus}
Let $X$ be a spectral space. Let $E \subset X$ be a constructible subset.
Let $W \subset X$ be the set of points of $X$ which specialize
to a point of $E$. Then $W \setminus E$ is a spectral space.
If $W = \bigcap U_i$ with $U_i$ as in
Lemma \ref{lemma-make-spectral-space}
(\ref{item-intersection-quasi-compact-open})
then $W \setminus E = \lim (U_i \setminus E)$.
\end{lemma}

\begin{proof}
Since $E$ is constructible, it is quasi-compact and hence
Lemma \ref{lemma-make-spectral-space} applies to $W$.
If $E$ is constructible, then $E$ is constructible in $U_i$
for all $i \in I$. Hence $U_i \setminus E$
is spectral by Lemma \ref{lemma-spectral-sub}.
Since $W \setminus E = \bigcap (U_i \setminus E)$ we have
$W \setminus E = \lim U_i \setminus E$ by the universal property of limits.
Then $W \setminus E$ is a spectral space by
Lemma \ref{lemma-directed-inverse-limit-spectral-spaces}.
\end{proof}

\end{document}
